\section{The potential layer and its boundary}
\label{sec:potential}

Pairing tries to assign one permanent response to every target.  A potential
argument instead assigns a weight to each surviving target and lets the
defender remove as much weight as possible on each placement.  The argument
is exact for a fixed finite family of windows.  On the infinite board, a
separate source term is needed when an attacker stone turns a stone-free
window into an active target.  This distinction determines which conclusions
are permanent and which cover only a finite horizon.

Throughout this section the defender is the Breaker in the potential
calculation.  A real defender completion only ends the game in the defender's
favor, so ignoring that objective is conservative.  We use
\[
 \lambda=\sqrt3.
\]
This is the largest exponential base supported by the two-against-two greedy
comparison of Erd\H{o}s--Selfridge and Beck
\cite{erdos-selfridge-1973,beck-2008}.

\subsection{Fixed target families}

Within this section, a window is \emph{potential-active} at $P$ when it is
attacker-touched and defender-free:
$W\cap A\ne\varnothing$ and $W\cap D=\varnothing$.  This is narrower than
``alive for $\A$'' in \cref{def:D5}; an all-empty window is not
potential-active.

The fixed-family potential retains window labels.  Distinct windows continue
to count separately even if their residual empty sets happen to coincide.
This makes the danger removed by a defender placement an exact sum.

\begin{definition}[Residual potential and danger]
\label{def:ES-potential}
Let $F$ be a finite labelled family of windows.  At a position $P$, let
\[
 \mathcal L_F(P)=\{W\in F:W\cap D=\varnothing\}
\]
be its surviving labels and put
\[
 E_P(W)=W\setminus A.
\]
The fixed-family residual potential is
\[
 \Psi_F(P)=\sum_{W\in\mathcal L_F(P)}
                   \lambda^{-|E_P(W)|}.
\tag{10.1}
\]
For an empty cell $x$, its fixed-family danger is
\[
 d_F(x;P)=
 \sum_{\substack{W\in\mathcal L_F(P)\\x\in E_P(W)}}
       \lambda^{-|E_P(W)|}.
\tag{10.2}
\]
The dynamic touched-window potential is
\[
 \Phi(P)=
 \sum_{\substack{W:W\cap A\ne\varnothing\\W\cap D=\varnothing}}
       \lambda^{-e_P(W)},
 \qquad e_P(W)=6-|W\cap A|.
\tag{10.3}
\]
Stone-free windows are absent from \textup{(10.3)}.
\end{definition}

\begin{plainterms}
Each unfinished target receives more weight as the attacker fills it.  A
defender move deletes every target it meets, while an attacker move increases
the terms of targets it advances.  The dynamic sum omits untouched windows
because the infinite board contains infinitely many of them; assigning each
even a fixed positive virgin weight would make the sum diverge.
\end{plainterms}

For a fixed residual family, a defender placement at $x$ decreases
$\Psi_F$ by exactly $d_F(x;P)$.  An attacker placement at $x$ multiplies each
affected term by $\lambda$, and therefore increases the potential by exactly
$(\lambda-1)d_F(x;P)$.  If all dangers are at most $S$ before an attacker
placement, every danger after that placement is at most $\lambda S$.

Sequential greedy defense recomputes \textup{(10.2)} before each defender
placement and takes an empty cell of maximum danger.  Ties follow one fixed
deterministic order.  If all dangers vanish, the strategy uses a deterministic
legal filler.  One concrete filler rule chooses an occupied cell with greatest
$q$ coordinate, breaks ties by greatest $r$, and takes its positive
$Q$-neighbor.  That neighbor is empty by maximality and is legal at distance
one.  The rule is recomputed if a second filler is needed.

The next lemma is the numerical core.  Its hypothesis about legal positive
danger cells holds when every surviving label was potential-active at the
initial defender epoch.

\begin{lemma}[One fixed-family round]
\label{lem:ES-round}
Let the defender make two sequential fixed-family greedy placements, followed
by the next attacker turn or its terminal prefix.  Suppose every
positive-danger cell is legal.  If the round reaches the next full defender
epoch, the fixed-family potential there is no larger than it was at the
initial defender epoch.  The corresponding upper bound also holds at a
terminal attacker prefix.
\end{lemma}

\begin{plainterms}
The two greedy defender placements remove at least as much potential as the
next attacker pair can restore.  Recomputing after the first defender move is
part of the argument, because it supplies two separate maxima.  The equality
$\lambda^2-1=2$ is the numerical balance behind the result.
\end{plainterms}

\begin{proof}
Let $\delta_1$ and $\delta_2$ be the exact reductions on the two defender
placements.  After both placements, let $S$ be the greatest remaining danger.
Deleting labels cannot increase danger at another cell.  Each greedy maximum
was therefore at least $S$, so
\[
 \delta_1+\delta_2\geq2S.
\]
The first attacker placement adds at most $(\lambda-1)S$.  The second adds at
most $(\lambda-1)\lambda S$.  Their total increase is at most
\[
 (\lambda-1)(1+\lambda)S=(\lambda^2-1)S=2S.
\]
The net change is nonpositive.  Omitting a terminal second placement can only
reduce the increase.
\end{proof}

The strict terminal threshold is one: a completed labelled target has empty
residual set and contributes $\lambda^0=1$ by itself.

\begin{theorem}[Fixed-family blocking]
\label{thm:ES-fixed-family}
Let $P_0$ be a nonterminal defender-\emph{FirstStone} position, and let $F$ be
a finite fixed family of windows that are potential-active at $P_0$.  If
\[
 \Psi_F(P_0)<1,
\tag{10.4}
\]
then sequential fixed-family greedy defense prevents the attacker from ever
completing a member of $F$.
\end{theorem}

\begin{plainterms}
The theorem protects exactly the targets named at the starting position.  It
does not automatically enroll a window that is still stone-free there.  The
full defender turn and the strict inequality are both part of the statement.
Fillers are permitted once no positive fixed-family danger remains.
\end{plainterms}

\begin{proof}
By \cref{lem:ES-round}, induction over full defender epochs keeps
$\Psi_F\leq\Psi_F(P_0)<1$.  Defender placements only decrease the sum.  If a
member of $F$ were completed on the second attacker placement, the next-epoch
potential would be at least one, a contradiction.  If completion occurred on
the first placement, the one-placement version of the round estimate gives
\[
 \Psi_F\leq \Psi_F(P_t)-2S+(\lambda-1)S<1
\]
at that placement, again contradicting its completed term of weight one.
Every alleged win occurs at a finite placement, so induction over finite
prefixes proves the infinite conclusion.
\end{proof}

The base $\sqrt3$ is optimal within this comparison.  For a general
$\lambda>1$, the attacker pair can add $(\lambda^2-1)S$, while two greedy
defender placements guarantee only $2S$.  Thus $\lambda\leq\sqrt3$ is needed,
and the largest permitted base minimizes every unfinished-edge weight.  Base
$2$ is not a substitute: three separated count-four targets have total weight
$3\cdot2^{-2}=3/4<1$, but the defender kills at most two of them.  The
inequality in \textup{(10.4)} also cannot be relaxed to equality at the correct
base.  Three pairwise separated count-four targets have total weight
$3\lambda^{-2}=1$ and give the same loss.

The theorem is sufficient, not necessary.  A single defender placement can
kill many labelled targets sharing its cell even when their total weight is at
least one.  The full-turn premise also matters.  With only one defender
placement remaining, two separated count-four targets have
$\Psi_F=2/3<1$; the defender kills one and the attacker completes the other.

\subsection{Birth mass and permanent local certificates}

The family in \cref{thm:ES-fixed-family} is fixed.  Dynamic touched-window
greedy enrolls targets during play, and their first contribution is not paid
by the fixed-family round estimate.  The exact missing quantity is their birth
mass.

Let $P_t$ be the start of a full defender turn.  During the following attacker
pair, let $N_t$ contain the windows that were stone-free at $P_t$, survived the
 intervening defender pair, and first became potential-active during that attacker
pair.  For $W\in N_t$, let $a_t(W)\in\{1,2\}$ be the number of stones from the
pair that lie in $W$ at the end of the pair or terminal prefix.  Define
\[
 C_t=\sum_{W\in N_t}\lambda^{-(6-a_t(W))}.
\tag{10.5}
\]
If $u_{1,t}$ windows are born with one stone and $u_{2,t}$ with two, then
\[
 C_t=\frac{u_{1,t}}{9\sqrt3}+\frac{u_{2,t}}9.
\]
The initial potential-active cohort obeys the fixed-family round estimate, while the new
cohort contributes exactly \textup{(10.5)}.  Hence every completed round
satisfies
\[
 \Phi(P_{t+1})\leq\Phi(P_t)+C_t.
\tag{10.6}
\]

\begin{theorem}[Cumulative-birth certificate]
\label{thm:ES-cumulative}
Let $P_0$ be a nonterminal defender-\emph{FirstStone} position.  Fix the
tie-breaking and filler policy above and use sequential dynamic touched-window
greedy from $P_0$.  Suppose there is a constant $B_\infty$ such that, uniformly
over every legal attacker strategy $\tau$,
\[
 \sup_\tau\sum_{t=0}^{\infty}C_t(\tau)\leq B_\infty,
 \qquad \Phi(P_0)+B_\infty<1.
\tag{10.7}
\]
Then the defender prevents every attacker completion.  More generally, if
\[
 \sup_\tau\sum_{t=0}^{T-1}C_t(\tau)\leq B_T,
 \qquad \Phi(P_0)+B_T<1,
\tag{10.8}
\]
then the first $T$ attacker pairs are safe.
\end{theorem}

\begin{plainterms}
New targets may be enrolled, but their total initial weight must be reserved
in advance.  The bound is uniform over attacker strategies; a small birth sum
observed on one line is not a certificate.  The theorem is useful when a
finite region or another argument supplies a finite activation reservoir.
\end{plainterms}

\begin{proof}
Before attacker pair $t$, repeated use of \textup{(10.6)} gives
\[
 \Phi(P_t)\leq\Phi(P_0)+\sum_{i=0}^{t-1}C_i<1.
\]
Apply the one-cycle result proved in \cref{thm:ES-one-cycle} below.  It makes
pair $t$ safe, after which \textup{(10.6)} advances the bound.  Induction proves
\textup{(10.8)}.  Under \textup{(10.7)}, the same induction applies to every
finite prefix produced against every $\tau$.
\end{proof}

A finite confinement gives a concrete reserve.  Fix a nonterminal
defender-\emph{FirstStone} position $P_0$ and a finite set $R$ of cells empty
at $P_0$.  Suppose every future attacker placement lies in $R$.  Put
\[
\begin{aligned}
 \mathcal A_R&=\{W:W\cap D_0=\varnothing,
   \ W\cap A_0\ne\varnothing,\ W\setminus A_0\subseteq R\},\\
 \mathcal V_R&=\{W:W\text{ is stone-free at }P_0,\ W\subseteq R\},
\end{aligned}
\tag{10.9}
\]
and let
\[
 \Phi_R(P_0)=\sum_{W\in\mathcal A_R}
       \lambda^{-e_{P_0}(W)},
 \qquad N_R=|\mathcal V_R|.
\]
For a matching $M$ of vertex-disjoint two-cell subsets of $R$, let $q(M)$ be
the number of windows in $\mathcal V_R$ containing both cells of some pair in
$M$, and define
\[
 P_R=\max_M q(M),\qquad
 B_R^*=N_R\lambda^{-5}
       +P_R(\lambda^{-4}-\lambda^{-5}).
\tag{10.10}
\]

\begin{theorem}[Finite-region certificate]
\label{thm:ES-finite-region}
Assume the attacker is forever confined to the finite set $R$.  If
\[
 \Phi_R(P_0)+B_R^*<1,
\tag{10.11}
\]
then delayed-enrollment region-target greedy prevents every attacker win.  The
weaker condition
\[
 \Phi_R(P_0)+\frac{N_R}{9}<1
\tag{10.12}
\]
is also sufficient.  Every positive-reduction defender placement is an empty
 cell of a currently potential-active member of
$\mathcal A_R\cup\mathcal V_R$ and is legal.  Fillers need not lie in $R$.
\end{theorem}

\begin{plainterms}
Confinement converts infinitely many possible virgin targets into a finite
list.  Each such target pays a one-stone birth charge, with an extra charge if
both of its first attacker stones arrive in one turn.  The theorem does not
prove confinement; a solver must verify that premise by another certificate.
\end{plainterms}

\begin{proof}[Proof sketch]
Every virgin target is enrolled once, with charge $\lambda^{-5}$, or
$\lambda^{-4}$ when both first stones arrive in one attacker turn.
Those exceptional birth pairs form a matching in $R$, so at most $P_R$ targets
receive the larger charge and the total is at most
$B_R^*\leq N_R\lambda^{-4}=N_R/9$.
For the potential $\Xi_R$ of the currently potential-active members of
$\mathcal A_R\cup\mathcal V_R$, delayed greedy gives
$\Xi_R(P_{t+1})\leq\Xi_R(P_t)+C_t^R$ and the charge gives
$\sum_tC_t^R\leq B_R^*$.
The induction of \cref{thm:ES-cumulative} therefore keeps $\Xi_R<1$.
Confinement places every possible completed window in this target family, and
every positive-danger cell is a legal empty within distance five of an
attacker stone.
The complete charging argument is in the companion
\cite{hexobot-companions}.
\end{proof}

For checking \textup{(10.11)}, $P_R$ must be exact or replaced by a proved
upper bound.  The elementary bound
\[
 P_R\leq\min\{N_R,5\lfloor |R|/2\rfloor\}
\]
follows because two distinct cells lie together in at most five length-six
windows.  One exhibited matching gives only a lower bound on $P_R$ and cannot
be inserted into \textup{(10.10)} as a certificate.

\subsection{Finite-horizon consequences}

Even without a birth reserve, a stone-free length-six window cannot be filled
in one attacker pair.  This gives the immediate dynamic result.

\begin{theorem}[One-cycle certificate]
\label{thm:ES-one-cycle}
If the defender is at \emph{FirstStone} and $\Phi(P_t)<1$, then two sequential
dynamic touched-window greedy placements prevent an attacker win during the
immediately following two-placement turn, including its first-placement win
check.
\end{theorem}

\begin{plainterms}
The current active targets are below the fixed-family threshold, so greedy
blocks them through the next attacker turn.  A target that is still
stone-free cannot acquire six stones in two placements.  The potential may
exceed one after those placements, so this is not a permanent certificate.
\end{plainterms}

\begin{proof}
Apply \cref{thm:ES-fixed-family} to the potential-active windows at $P_t$.  A window
outside that cohort either contains a permanent defender stone or is
stone-free at $P_t$.  The latter receives at most two attacker stones during
the turn and cannot be completed.
\end{proof}

Fixing the initial cohort gives a longer horizon because a virgin target
needs six attacker placements, regardless of where the placements are split
into turns.

\begin{theorem}[Five-placement certificate]
\label{thm:ES-five-placement}
If the defender is at \emph{FirstStone} and $\Phi(P_0)<1$, then the defender
has a strategy preventing an attacker win during the next five attacker
placements: two complete attacker turns and the first placement of the third.
\end{theorem}

\begin{plainterms}
The defender uses fixed-family greedy on the targets active at the start.
Those targets stay blocked forever.  Every other live candidate began empty
and still lacks at least one attacker stone after five future placements.
The sixth placement is outside this guarantee.
\end{plainterms}

\begin{proof}
Let $F$ be the potential-active family at $P_0$ and use fixed-$F$ sequential greedy.
The initial potential of $F$ is $\Phi(P_0)<1$, so
\cref{thm:ES-fixed-family} blocks it forever.  A window containing a defender
stone at $P_0$ remains dead.  Every remaining window was stone-free and has at
most five attacker stones after five future attacker placements.
\end{proof}

The three-pair result uses a deliberate defender placement against virgin
targets and only one greedy placement on the same turn.  The following lemma
gives the exact cost of that split.

\begin{lemma}[One-greedy amplification]
\label{lem:ES-one-greedy}
Let $F$ be a fixed potential-active family at a full defender epoch and write
$X=\Psi_F$.  Suppose a defender turn consists of one arbitrary legal placement
followed by one $F$-greedy placement, and then at most two attacker placements.
At every attacker prefix, including a terminal placement,
\[
 \Psi_F\leq\frac32X.
\tag{10.13}
\]
If both attacker placements are nonterminal, this is also the bound at the
next defender epoch.  If $3X/2<1$, no member of $F$ is completed at either
placement.
\end{lemma}

\begin{plainterms}
One defender placement may be spent on a virgin target rather than on
fixed-family danger.  The remaining greedy move still limits the fixed-family
growth to a factor of three halves for that round.  This controlled loss of
margin permits the three-pair strategy to mix geometric blocking with
potential defense.
\end{plainterms}

\begin{proof}
The arbitrary defender placement leaves a potential $X'\leq X$.  Let the
greedy placement reduce it by $\delta$, and let $S$ be the greatest remaining
danger.  Deletion monotonicity and greediness give $S\leq\delta$.  A cell's
danger is also a subsum of the remaining potential, so
$S\leq X'-\delta$.  The next attacker pair adds at most $2S$ by the same
calculation as \cref{lem:ES-round}.  Therefore
\[
 \Psi_F^+
 \leq X'-\delta+2\min\{\delta,X'-\delta\}
 \leq\frac32X'\leq\frac32X.
\]
After only the first attacker placement, its increase is at most
$(\lambda-1)S\leq2S$, so the same bound holds.  A completed member contributes
one and contradicts \textup{(10.13)} when $3X/2<1$.
\end{proof}

The geometry of the first attacker pair identifies every virgin window that
could be completed on the sixth future attacker placement.  Blocking those
windows costs zero, one, or two split defender turns according to the pair's
axis relation.

\begin{theorem}[Three-pair certificate]
\label{thm:ES-three-pair}
Let $P_0$ be a nonterminal defender-\emph{FirstStone} position, and let $x,y$
be the first future attacker pair.  The defender can prevent a win during the
next three complete attacker turns under the respective conditions
\[
 \Phi(P_0)<
 \begin{cases}
  1, & x,y\text{ lie in no common window},\\
  2/3, & x,y\text{ have common-window axis distance }2,3,4,\text{ or }5,\\
  4/9, & x,y\text{ are axis-adjacent}.
 \end{cases}
\tag{10.14}
\]
In particular, $\Phi(P_0)<4/9$ unconditionally certifies all three attacker
turns.
\end{theorem}

\begin{plainterms}
A virgin window can first win on the sixth future attacker placement.  Its
first two attacker cells reveal which virgin windows could contain all six
placements.  The defender blocks that small common family while retaining
enough fixed-family potential margin.  Adjacent first cells expose five
possible windows and require two split defender turns, which accounts for the
smaller threshold $4/9$.
\end{plainterms}

\begin{proof}[Proof sketch]
Let $F$ be the initial potential-active family and $X_t=\Psi_F(P_t)$; the first greedy
defender pair gives $X_1\leq X_0$, while a virgin sixth-placement completion
must contain both $x$ and $y$.
If they share no window, later $F$-greedy pairs suffice.
At common-window distance $2$ through $5$, write $y=x+dv$ for a signed unit
axis vector $v$; playing $x+v$ and then one greedy cell hits every eligible
virgin target and gives $X_2\leq3X_1/2<1$.
When $x,y$ are adjacent, playing $x-v$ leaves at most
$W_*=\{x,x+v,\ldots,x+5v\}$, which a later deliberate hit blocks, and two
one-greedy estimates give $X_3\leq(3/2)^2X_1\leq9X_0/4<1$.
Thus the initial family remains below threshold and every eligible virgin
target is hit.
The companion gives the legality and filler cases in full
\cite{hexobot-companions}.
\end{proof}

\subsection{How narrow is the raw threshold?}

For $s\in\{1,\ldots,5\}$ let $n_s$ count potential-active windows with exactly
$s$ attacker stones.  Then
\[
 \Phi=\sum_{s=1}^{5}n_s\lambda^{s-6},
\]
with the weights shown in \cref{tab:potential-weights}.

\begin{table}[tb]
\centering
\caption{One-window contributions to the dynamic potential.}
\label{tab:potential-weights}
\begin{tabular}{cccc}
\toprule
attacker stones & empties & exact weight & decimal\\
\midrule
$1$ & $5$ & $1/(9\sqrt3)$ & $0.064150$\\
$2$ & $4$ & $1/9$ & $0.111111$\\
$3$ & $3$ & $1/(3\sqrt3)$ & $0.192450$\\
$4$ & $2$ & $1/3$ & $0.333333$\\
$5$ & $1$ & $1/\sqrt3$ & $0.577350$\\
\bottomrule
\end{tabular}
\end{table}

If $I=\sum_{s=1}^{5}s n_s$ is the number of live
attacker-stone--window incidences, direct substitution gives
$\lambda^{s-6}/s\geq1/18$, with equality at $s=2$.  Thus
\[
 \Phi\geq I/18,
\]
and $\Phi<1$ forces $I\leq17$.  Since each attacker stone has eighteen window
incidences, the raw threshold occurs only after almost all such incidences have
been killed by defender stones.  It is not a generic opening or middlegame
condition.

The comparison can be made without floating point.  Set
\[
 a=n_1+3n_3+9n_5,
 \qquad b=n_2+3n_4.
\]
Then $\Phi<1$ holds exactly when
\[
 b\leq8\qquad\text{and}\qquad a^2<3(9-b)^2.
\tag{10.15}
\]
Indeed, after multiplication by $9\sqrt3$, the inequality is
$a<\sqrt3(9-b)$.  Its right side must be positive, after which squaring
preserves the comparison.

\subsection{Dynamic greedy can be forced to lose}

The fixed-family theorem does not prove that repeatedly recomputing the
dynamic touched family is a global defense.  The next finite counterexample
branches over every exact greedy maximum.  All maxima are positive, so the
result is independent of the filler policy.

\begin{theorem}[All-ties greedy refutation]
\label{thm:ES-greedy-refutation}
Let
\[
 A_0=\{(0,0)\},\qquad D_0=\{(1,0)\},
\]
with the defender at \emph{FirstStone}.  Then $\Phi(P_0)<1$.  Against every
sequential dynamic touched-window greedy defense, the fixed attacker
continuation
\[
\begin{array}{c@{\qquad}l}
\text{turn }1 &(2,-4),(2,2),\\
\text{turn }2 &(-5,0),(-4,0),\\
\text{turn }3 &(-3,0),(-2,0),\\
\text{turn }4 &(-1,0)
\end{array}
\tag{10.16}
\]
wins on the first placement of turn $4$.  The completed window is
$\{(-5,0),(-4,0),\ldots,(0,0)\}$.
\end{theorem}

\begin{plainterms}
Greedy is not rescued by choosing ties more carefully.  The verifier branches
over every maximum-danger reply and reaches the same fixed winning placement
on all 124 surviving boards.  This refutes the dynamic greedy strategy, not
the claim that some different defender strategy may survive from every
position with $\Phi<1$.
\end{plainterms}

\begin{proof}[Machine-verified finite proof]
The defender stone $(1,0)$ kills five of the eighteen windows through
$(0,0)$.  The initial profile is therefore
$(n_1,n_2,n_3,n_4,n_5)=(13,0,0,0,0)$, and
\[
 \Phi(P_0)=\frac{13\sqrt3}{27}
           =\frac{13}{9\sqrt3}<1,
\]
because $13^2=169<243=3\cdot9^2$.  Every attacker placement in
\textup{(10.16)} is legal on every branch.  The first two are at distance four
from $(0,0)$, and each later placement is within distance five of an attacker
stone on the target line.

The verifier represents every potential and danger as an integer pair
$(a,b)$ denoting $(a+b\sqrt3)/27$.  Comparisons use integer arithmetic.  When
the coefficients of a difference have opposite signs, the verifier compares
$a^2$ with $3b^2$; irrationality excludes equality.  It enumerates every
maximum at each defender placement, merges only identical colored boards at
the same phase, and then applies the fixed attacker continuation.  The exact
state counts and maxima are shown in \cref{tab:greedy-branches}.  There is no
random choice, beam limit, or depth cutoff.

\begin{table}[tb]
\centering
\caption{Exhaustive dynamic-greedy branch counts.  Danger values are scaled by
$27$.}
\label{tab:greedy-branches}
\begin{tabular}{crlr}
\toprule
defender placement & input boards & exact maximum distribution & output boards\\
\midrule
$D0.1$ & $1$ & $5\sqrt3$ & $4$\\
$D0.2$ & $4$ & $5\sqrt3\ (4)$ & $4$\\
$D1.1$ & $4$ & $6\sqrt3\ (2),\ 7\sqrt3\ (2)$ & $12$\\
$D1.2$ & $12$ & $5\sqrt3\ (12)$ & $124$\\
$D2.1$ & $124$ & $10\sqrt3\ (124)$ & $242$\\
$D2.2$ & $242$ & $9\sqrt3\ (6),\ 10\sqrt3\ (236)$ & $124$\\
$D3.1$ & $124$ & $21+4\sqrt3\ (124)$ & $124$\\
$D3.2$ & $124$ & $3+9\sqrt3\ (106),\ 10\sqrt3\ (18)$ & $124$\\
\bottomrule
\end{tabular}
\end{table}

The last defender turn contains the trap.  Immediately before $D3.1$, the
cell $(-6,0)$ lies in five potential-active horizontal windows with attacker counts
$1,2,3,4,4$.  Hence
\[
 27d(-6,0)=\sqrt3+3+3\sqrt3+9+9=21+4\sqrt3.
\tag{10.17}
\]
The winning cell $(-1,0)$ lies in one count-four window and the count-five
target, so
\[
 27d(-1,0)=9+9\sqrt3.
\tag{10.18}
\]
Their difference is $12-5\sqrt3>0$, since $12^2>3\cdot5^2$.  Thus every
branch chooses $(-6,0)$ uniquely.  That stone kills the count-four term at
$(-1,0)$ and leaves
\[
 27d(-1,0)=9\sqrt3.
\]
On 106 boards the unique second maximum is $(-3,1)$ with danger
$3+9\sqrt3$.  On the other 18 it is $(-2,-1)$ with danger $10\sqrt3$.
Both exceed $9\sqrt3$.  No greedy branch occupies $(-1,0)$.

\begin{figure}[tb]
\centering
\begin{tikzpicture}[scale=.48]
  \HexRhombus{-6}{2}{-4}{2}
  \HexWindow[window one]{-5}{0}{1}{0}
  \HexAttackerList{-5/0,-4/0,-3/0,-2/0,0/0,2/-4,2/2}
  \HexMark{-6}{0}{cross}{}
  \HexMark{-1}{0}{cross}{}
  \node[font=\small,align=center,above=4mm] at \HexAt{-6}{0}
    {$27d=21+4\sqrt3$\\greedy maximum};
  \node[font=\small,align=center,below=4mm] at \HexAt{-1}{0}
    {$27d=9+9\sqrt3$\\winning cell};
  \node[font=\small,anchor=west] at \HexAt{-2}{.85}
    {target $W$};
\end{tikzpicture}
\caption{The common attacker geometry on all 124 boards immediately before
$D3.1$.  The attacker set is
$\{(-5,0),(-4,0),(-3,0),(-2,0),(0,0),(2,-4),(2,2)\}$.  Defender stones are
omitted because their locations depend on the greedy branch.  The target,
candidate cells, and displayed danger values are the same on every branch.}
\label{fig:greedy-trap}
\end{figure}

The verifier also checks radius-eight legality, the absence of every attacker
or defender completion before the terminal placement, every potential-active-window
profile, every exact maximizer, and the terminal window.  At the terminal
placement all 124 boards have exactly one completed attacker window and hence
$\Phi\geq1$.  The scheduled attacker cell is absent from every prior
maximizer union.  This exhausts every dynamic-greedy tie-breaking sequence and
proves the theorem.  The exact checker and its full profile trace are included
with the companion proofs \cite{hexobot-companions}.
\end{proof}

The compact position has unequal material counts, but the same tree occurs in
a legal engine history.

\begin{proposition}[Reachable greedy witness]
\label{prop:ES-reachable-greedy}
There is a legal $39$-placement history ending at a nonterminal defender-
\emph{FirstStone} position with $20$ attacker stones, $19$ defender stones,
the same initial profile $(13,0,0,0,0)$, and exactly a translated copy of the
greedy tree in \cref{thm:ES-greedy-refutation}.
\end{proposition}

\begin{proof}[Proof sketch]
Translate the compact core by $(-1,0)$, so its defender stone is the opening
at $(0,0)$ and its attacker stone is $(-1,0)$.  Let
\[
 c=(-11,10),\quad I=B_2(c),\quad
 R=B_3(c)\setminus B_2(c),\quad g=(-9,8).
\]
Here $|I|=19$ and $|R|=18$; after the opening, the attacker takes $(-1,0)$ and
$g$, then nine cycles alternate two unused cells of $R$ with two unused cells
of $I\setminus\{g\}$.
All moves are legal because $d(g,(-1,0))=8$ and every cell of $B_3(c)$ is
within five of $g$.
No prefix has a six-run, and the final defender ring blocks every window
meeting $I$ while distance-six padding isolates the translated tactical core.
The count is $1+2+9\cdot4=39$ placements with material totals $20$ and $19$.
The companion checker verifies every prefix and reruns the full greedy tree
\cite{hexobot-companions}.
\end{proof}

\subsection{Clean escape and failure of raw renewal}

The greedy counterexample concerns one strategy.  A separate geometric lemma
shows that no strategy can keep the raw sublevel set $\Phi<1$ invariant at
every defender epoch.  It still does not give an attacker win against every
non-greedy defense.

\begin{lemma}[Universal clean escape]
\label{lem:ES-clean-escape}
From every finite, nonempty, nonterminal position immediately before an
attacker turn, the attacker has a legal pair that births exactly $36$ distinct
one-stone windows.  Its birth mass is
\[
 36\lambda^{-5}=\frac4{\sqrt3}.
\tag{10.19}
\]
\end{lemma}

\begin{plainterms}
The attacker can move twice toward an unoccupied extreme of the finite board.
Both local stars are clean and share no target window.  The resulting birth
mass already exceeds one, so raw current potential cannot renew itself as an
epoch-by-epoch invariant under any defender policy.
\end{plainterms}

\begin{proof}
Put $h(q,r)=q+r$.  Choose an occupied cell $z$ maximizing $h$, and write
$M=h(z)$.  Play
\[
 x=z+(4,4),\qquad y=x+(4,4).
\]
Both displacements have hex distance eight, so the placements are legal in
sequence.  They are empty because their $h$-values exceed $M$.

A horizontal or vertical length-six window through $x$ has minimum $h$ at
least $h(x)-5=M+3$.  A diagonal window through $x$ has constant
$h=M+8$.  Every such window was therefore stone-free.  A horizontal or
vertical window through $y$ has minimum $h$ at least $M+11$, while a diagonal
window through $y$ has constant $h=M+16$.  Those windows were stone-free even
after the placement at $x$.  Finally, $y-x=(4,4)$ lies on none of the three
axes, so $x$ and $y$ share no window.  Each placement births its eighteen
incident windows, all with one attacker stone.  The $36$ distinct terms have
total weight
\[
 36(\sqrt3)^{-5}=\frac4{\sqrt3}.
\]
\end{proof}

Starting from one defender stone and no attacker stone, $\Phi=0$.  After any
defender pair, \cref{lem:ES-clean-escape} makes the next defender-epoch
potential at least $4/\sqrt3>1$.  Repeating the construction on every
continuing attacker turn creates 36 fresh labels each time.  The cumulative
birth sum diverges, even though each such label may remain at count one or be
killed.  Radius-eight mobility alone therefore supplies no summable reserve
for \cref{thm:ES-cumulative}.

\subsection{Two static repair attempts}

The first no-go result strengthens \cref{thm:T8}.  Finite occupation cannot
remove enough exceptional windows to make a static matching possible on the
remaining virgin board.

\begin{proposition}[Finite saturation does not enable static pairing]
\label{prop:ES-no-static-pairing}
No partial matching covers all but finitely many length-six windows.  Hence no
finite initial position admits a static pairing that covers every initially
virgin target.
\end{proposition}

\begin{proof}
Suppose a matching misses at most $K$ windows.  Discard its pairs that lie in
no length-six window.  Let $B_R$ be the radius-$R$ hex ball, with
\[
 |B_R|=3R^2+3R+1.
\]
For each axis and $R\geq5$, the $2R+1$ line segments through $B_R$ contain
\[
 |B_R|-5(2R+1)
\]
length-six windows wholly inside the ball.  Across three axes, at least
\[
 3|B_R|-15(2R+1)-K
\]
of these internal windows must be covered.  One matched pair is contained in
at most five windows on its axis.  The ball therefore contains at least
\[
 \frac25\bigl(3|B_R|-15(2R+1)-K\bigr)
\]
matched-cell incidences.  A matching supplies at most $|B_R|$, so
\[
 |B_R|\leq60R+30+2K.
\]
This contradicts $|B_R|=3R^2+3R+1$ for sufficiently large $R$.

If a finite position has occupied support $S$, at most $18|S|$ windows meet
that support.  Every other window is initially virgin.  A matching covering
all virgin windows would miss only finitely many windows, contrary to the
first part.
\end{proof}

The second no-go result concerns a static additive account that tries to put
a summable, spatially damped weight on every virgin window while retaining the
same factor-three two-placement growth comparison.

\begin{proposition}[Static two-tier damping no-go]
\label{prop:ES-no-static-damping}
Suppose a nonnegative static per-window account assigns weight $w_s(W)$ when
$W$ contains $s$ attacker stones.  Assume
\[
 \sum_W w_0(W)<\infty,\qquad w_6(W)\geq1,
\]
and, for every window,
\[
 w_{s+2}(W)\leq3w_s(W),\qquad s=0,2,4.
\tag{10.20}
\]
No such account exists.
\end{proposition}

\begin{proof}
For every window, the terminal bound and the three growth inequalities give
\[
 1\leq w_6(W)\leq3w_4(W)\leq9w_2(W)\leq27w_0(W).
\]
Thus $w_0(W)\geq1/27$ for each of infinitely many windows.  Their virgin sum
cannot be finite.
\end{proof}

The proposition leaves other accounting systems open.  Its hypotheses do not
cover age-dependent weights, moving centers, cancellation between windows, a
nonadditive account, or an adaptive pairing whose ownership changes during
play.

\subsection{Finite reduction and the remaining boundary}

Although the unrestricted question has no known scalar renewal, every fixed
horizon occupies a finite part of the ideal board.  Compactness then relates
all finite survival policies to one infinite policy.

\begin{theorem}[Finite-horizon equivalence]
\label{thm:ES-finite-reduction}
Fix a finite nonempty position $P$ with occupied support $S_0$ of size $n$, at
defender-\emph{FirstStone}.  For $h\geq1$, define
\[
 R_h=\{x:\dist(x,S_0)\leq32h\}.
\tag{10.21}
\]
Every move in the first $h$ full defender/attacker rounds lies in $R_h$, and
\[
 |R_h|\leq n(3072h^2+96h+1).
\tag{10.22}
\]
The defender blocks forever from $P$ if and only if the defender has a
survival strategy for every finite $h$-round game on $R_h$.  The set $R_h$ is
a proved move universe through that horizon, not an added board boundary.
\end{theorem}

\begin{plainterms}
Every fixed-depth game can be searched exactly on a finite board region.  If a
survival policy exists at every depth, K\"onig's lemma selects policies whose
finite restrictions agree and joins them into one infinite strategy.  The
theorem gives no depth after which survival automatically implies permanent
safety; checking every horizon restates the infinite problem rather than
solving it algorithmically.
\end{plainterms}

\begin{proof}
The first $h$ rounds contain at most $4h$ placements.  Induct on the future
placement index $k$.  A legal placement lies within distance eight of an
occupied cell, so the $k$th placement lies within distance $8k$ of $S_0$.
Since $k\leq4h$, every move lies in $R_h$.

A radius-$r$ hex ball contains $3r^2+3r+1$ cells.  The union bound over the
$n$ initial support cells with $r=32h$ gives
\[
 |R_h|\leq n\bigl(3(32h)^2+3(32h)+1\bigr)
 =n(3072h^2+96h+1).
\]
The truncated legal game tree is finite, so backward induction decides it.

One implication is immediate: a forever-surviving strategy restricts to a
survival strategy for each horizon.  For the converse, form a rooted tree of
finite defender policy tables.  Level $h$ contains the policies that survive
the $h$-round game, and an edge joins a policy to its restriction at level
$h-1$.  Every level is finite, every node has finitely many extensions, and
every level is nonempty by hypothesis.  K\"onig's lemma supplies an infinite
chain of compatible policies.  Their union is one defender strategy that
survives every finite prefix and hence blocks forever.  Conversely, if no
forever strategy exists, some finite level is empty; finite minimax then gives
an attacker strategy forcing completion within that horizon.
\end{proof}

The results above leave five separate proof obligations.  They should not be
collapsed into one generic statement that the potential method is incomplete.

\paragraph{GAP-RAW.}
\Cref{thm:ES-greedy-refutation} defeats every dynamic touched-window greedy
tie-breaking, but it gives no attacker strategy against arbitrary non-greedy
play.  No theorem here proves that a non-greedy forever defense always exists.
The raw existential statement remains open.

\paragraph{GAP-GLOBAL-RENEWAL.}
\Cref{thm:ES-five-placement,thm:ES-three-pair} extend the certified horizon,
but neither returns to a defender epoch satisfying the same raw hypothesis.
\Cref{lem:ES-clean-escape} proves that no defense can use $\Phi<1$ itself as
the renewed invariant, even when the starting value is zero.

\paragraph{GAP-AMORTIZED-ABANDONMENT.}
An online account must discount indefinitely many weak births that the
attacker abandons, yet remain sound if dormant stones later participate in
multi-axis forks.  Clean escape produces a positive exact birth debit with no
maturity promotion.  \Cref{prop:ES-no-static-damping} excludes only the stated
static factor-three accounts and supplies no history-sensitive refund rule.

\paragraph{GAP-DYNAMIC-PAIRING.}
\Cref{prop:ES-no-static-pairing} excludes a static matching even after finite
initial exceptions.  It does not exclude an adaptive response system whose
pairs or ownership assignments change with the play.

\paragraph{GAP-FINITE-CUTOFF.}
\Cref{thm:ES-finite-reduction} proves that every loss has a finite witness and
that each fixed horizon is decidable.  It gives no computable or uniform bound
$H(P)$ such that survival through $H(P)$ rounds implies survival forever.

The strategy-level nonclaims also affect implementation.  A raw node-local
$\Phi<1$ test may certify \cref{thm:ES-one-cycle} or seed one of the stated
finite-horizon strategies.  It must not be converted into an unrestricted
defender win, draw, or global move-set pruning rule.  Every positive greedy
move lies in an empty cell of a currently potential-active monitored window and
is legal.  If no such target remains, one or both mandatory placements can be
fillers.  A solver implementing a potential-certified strategy must retain a
legal filler fallback, and it must account for the filler's effect on the
future radius-eight legality frontier.  Full derivations and the exact finite
traces appear in the companion proof documents \cite{hexobot-companions}.
